\section{Encoding Proof Obligations}
\label{sec:beer-proof-obligations}

Term soundness is the kernel; this section wraps it into the end-to-end guarantee.
A decoded proof obligation (\Cref{ch:b}) consists of a typing context, a list of definitions and hypotheses, and one or more goals.
The encoder {\small\leancodeptr{encode}{Encoder/Encoder.html\#encode}} processes a whole B environment as the composition
\[
  \mathsf{encodeTypeContext} \mathbin{;} \mathsf{encodeDefs} \mathbin{;} \mathsf{encodeDistinctFinite} \mathbin{;} \mathsf{encodeProofObligations} \mathbin{;} \mathsf{finalBulkDeclare},
\]
each stage emitting declarations or assertions into the growing script.

\subsection{From terms to environments}
\label{sec:beer-po-environments}

The typing context is emitted first: each variable \(\var{v}\) of B type \(\alpha\) yields \((\mathsf{declare\text{-}const}\ \var{v}\ \toSMT{\alpha})\), except that flagged variables (\Cref{sec:beer-rules-cast}) are declared at their tightened functional type.
Named definitions become \(\mathsf{define\text{-}fun}\)s over their encoded bodies, and the environment's distinctness and finiteness constraints are encoded and asserted.
Global \emph{hypotheses}, however, are deliberately not asserted at this stage: in the Atelier~B format, a clause that is a hypothesis for one proof obligation (an invariant, an assertion) is the \emph{goal} of another, so asserting them globally would beg the very questions the obligations pose.
Each proof obligation instead carries the list of definition blocks it may use, and only those are asserted, locally.
A final pass declares the residual variables that were typed during encoding but not yet declared---bound variables of constraints whose helper specifications were emitted globally.

\subsection{Goals and refutation}
\label{sec:beer-po-refutation}

Each proof obligation is encoded inside its own \((\mathsf{push}\ 1) \dots (\mathsf{pop}\ 1)\) block containing its local declarations, its asserted definitions and hypotheses, and then, per goal, a nested block asserting the \emph{negated} encoded goal followed by \(\mathsf{check\text{-}sat}\).
The scoping is not cosmetic.
Atelier~B synthesizes fresh identifiers per obligation and freely reuses a name at \emph{different types} across obligations; the local blocks realize the corresponding scoping discipline on the SMT side, so that each \(\mathsf{check\text{-}sat}\) sees exactly the declarations and assertions of its own obligation, and the typing invariants of \Cref{thm:beer-encodeTerm-spec} hold per block.
One script thus carries one \(\mathsf{check\text{-}sat}\) per goal, and each answer is interpreted independently: \textsf{unsat} discharges that goal, \textsf{sat} exhibits a countermodel and therefore invalidates it, while \textsf{unknown} or a timeout leaves it open.

On the running example, the emitted script is, abridged to its structure:
\begin{smtlisting}
(set-logic HO_ALL)
(declare-datatype Pair (par (T1 T2) ((pair (fst T1) (snd T2)))))
(declare-datatype Option (par (T) ((some (the T)) (none))))
(declare-const X (-> Int Bool))
(declare-const Y (-> Int Bool))
(push 1)
  (declare-const f (-> Int (Option Int)))
  (declare-const g (-> (Pair Int Int) Bool))
  (assert ...)  ; f : X --> Y, elided
  (assert ...)  ; g : X <-> Y, elided
  (push 1)
    (assert (not
      (forall ((z (Pair Int Int)))
      (forall ((f! (-> (Pair Int Int) Bool)))
        (=> (forall ((u Int) (v Int))
              (= (f! (pair u v)) (= (f u) (some v))))
        (=> ((lambda ((p (Pair Int Int))) (or (f! p) (g p))) z)
            (exists ((a Int) (b Int))
              (and (X a) (and (Y b) (= z (pair a b)))))))))))
    (check-sat)
  (pop 1)
(pop 1)
\end{smtlisting}
Every ingredient of the chapter is visible: the flagged \(\var{f}\) declared functionally and \(\var{g}\) relationally, the goal encoded as a guarded universal \cref{eq:beer-forall-shape} with the helper \(\lvar{f}\) universally bound under its graph specification, the union \cref{eq:beer-union-graph} applied to the bound variable, and the Cartesian product predicate of \Cref{sec:beer-rules-sets} as the conclusion.

\subsection{Semantic adequacy of proof-obligation encoding}
\label{sec:beer-po-soundness}

Validity of a proof obligation is measured against the B semantics.

\begin{definition}[Validity]
  \label{def:beer-po-validity}
  A proof obligation with typing context \(\Gamma_{\mathsf{B}}\), hypotheses \(H_1, \dots, H_n\), and goal \(G\) is \emph{valid} if for every valuation \(\Delta\) well formed for \(\Gamma_{\mathsf{B}}\) and defined on the free variables of the obligation,
  \[
    \Bigl(\textstyle\bigwedge_{i} \interpB{\aterm{H_i}{\Delta}} = \denval{\topZ}[\boolB]\Bigr)
    \implies
    \interpB{\aterm{G}{\Delta}} = \denval{\topZ}[\boolB].
  \]
  Hypotheses and goals are assumed well-defined, consistently with Atelier~B's separate emission of well-definedness obligations (\Cref{sec:b-wd}).
\end{definition}

\begin{theorem}[Semantic adequacy of proof-obligation encoding]
  \label{thm:beer-end-to-end}
  The refutation script emitted for a proof obligation is satisfiable with respect to the SMT semantics of \Cref{ch:smt} if and only if the proof obligation is invalid with respect to the B semantics.
  Equivalently, the emitted script is unsatisfiable if and only if the proof obligation is valid.
\end{theorem}

\begin{proof}[Proof sketch]
  For the forward implication, let an SMT valuation satisfy the emitted script.
  Transport the interpretations of the source variables back through the semantic bridge, and use the exactness specifications of the cast helpers to recover the corresponding B valuation.
  Semantic equivalence of term encoding then turns every true encoded hypothesis into a true B hypothesis and the asserted encoded negation of the goal into a false B goal.
  The obligation is therefore invalid.

  Conversely, invalidity yields a B valuation \(\Delta\) under which every hypothesis denotes \(\topZ\) and the goal denotes \(\botZ\).
  Transport \(\Delta\) through the bridge and interpret each helper by the semantic value determined by its specification: loosened variables receive the cast image of their source (\Cref{thm:loosening-correctness}), and function witnesses receive the functional collapse of the corresponding graph.
  Under the resulting SMT valuation, every encoded hypothesis denotes \(\topZ\), every helper specification is satisfied, and the encoded negation of the goal denotes \(\notZ\botZ = \topZ\).
  Hence all assertions in the obligation's block are satisfied and the emitted script is satisfiable.
\end{proof}

The theorem statement deserves the same honesty as the rest of the chapter.
The mechanized artifact is \Cref{thm:beer-encodeTerm-spec}, which carries, in its totality clause and helper-witness constructions, the term-level ingredients the equivalence consumes; its assembly at the level of scripts---iterating over assertions and interpreting \(\mathsf{push}/\mathsf{pop}\)---is a paper argument and is not, at the time of writing, a Lean theorem.
Moreover, satisfiability and unsatisfiability in \Cref{thm:beer-end-to-end} refer to our ZF semantics of SMT scripts, so interpreting the solver's answer assumes the solver correct for that semantics.
Both points are recorded in the trusted-base audit of \Cref{sec:beer-tcb}.

On the running example, the goal of \cref{eq:beer-running-po} is indeed valid---\(\interpB{\var{f}} \subseteq X \timesZ Y\) and \(\interpB{\var{g}} \subseteq X \timesZ Y\) follow from the two hypotheses, and the union of two subsets of \(X \timesZ Y\) is one---and \textsf{cvc5} answers \textsf{unsat} on the emitted script, closing the proof obligation.
